% Include with % Include with % Include with % Include with \input{zero_shot_datasets}; requires booktabs in the preamble.
% Bibliography: zero_shot_datasets.bib.
% Draft methods/protocol text, not a report of completed inference experiments.
% Verify all protocol claims against the actual runs before paper submission.
\subsection{Datasets for zero-shot cross-dataset evaluation}
\label{sec:zero-shot-datasets}

We prepare two target datasets for evaluating cross-dataset generalization:
ACDC~\cite{sakaridis2021acdc,sakaridis2026acdc} and
MRTMD~\cite{bugeja2025mrtmd}. Both are represented using the detector's six
classes, in output-index order: person, bicycle, motorcycle, car, bus, and truck.
Here, zero-shot denotes evaluation on a target dataset without target-specific
fitting or hyperparameter selection; it does not denote recognition of unseen
object categories. The proposed evaluation protocol reserves these datasets
for evaluation and fixes preprocessing, checkpoint selection, and detection
thresholds before examining target-set results.

\paragraph{ACDC.}
ACDC comprises driving scenes captured in Switzerland under adverse visual
conditions. We prepare its complete adverse-condition validation split of
406 images at $1920\times1080$ pixels: 100 fog, 106 night, 100 rain, and 100
snow images. Normal-condition reference images are not included.
The source person and rider categories are mapped to person, with bicycle and
motorcycle boxes retained separately. The 46 train annotations are removed
while their images are retained. The resulting six-class annotation set
contains 2,685 individual objects and 63 crowd regions. Native JSON annotations
retain the source boxes, areas, and crowd flags. This mapping defines a custom
six-class evaluation and is not directly comparable with ACDC's official
eight-class detection benchmark.

\paragraph{MRTMD.}
MRTMD comprises traffic-video frames curated from YouTube and distributed at
multiple resolutions. We prepare only the labelled 1080p release, avoiding
repeated counting of different-resolution copies of the same frames.
The pinned source annotations contain 3,733 image records and 69,348 boxes.
Removing 22 redundant image records with identical annotations yields 3,711
images and 69,010 boxes. Class IDs are reordered without merging or dropping
categories. We clip 575 boxes with boundary overflows of at most 0.25 pixels,
and recompute areas from the corrected 1080p boxes because the published area
fields are approximately four times their 1080p bounding-box areas.
The original annotations are retained separately.

\begin{table}[t]
  \centering
  \small
  \caption{Prepared evaluation sets. Crowd annotations describe regions and
  are reported separately from individual object boxes.}
  \label{tab:zero-shot-dataset-sizes}
  \begin{tabular}{lrrr}
    \toprule
    Dataset & Images & Objects & Crowds \\
    \midrule
    ACDC validation & 406 & 2,685 & 63 \\
    MRTMD 1080p & 3,711 & 69,010 & 0 \\
    \bottomrule
  \end{tabular}
\end{table}

\begin{table}[t]
  \centering
  \small
  \caption{Class composition after mapping. ACDC object and crowd counts are
  separate; all MRTMD annotations are individual object boxes.}
  \label{tab:zero-shot-class-counts}
  \begin{tabular}{lrrr}
    \toprule
    Class & \shortstack{ACDC\\objects} & \shortstack{ACDC\\crowds} & MRTMD \\
    \midrule
    Person & 567 & 11 & 21,274 \\
    Bicycle & 68 & 22 & 64 \\
    Motorcycle & 59 & 1 & 15,525 \\
    Car & 1,836 & 29 & 30,187 \\
    Bus & 49 & 0 & 337 \\
    Truck & 106 & 0 & 1,623 \\
    \midrule
    Total & 2,685 & 63 & 69,010 \\
    \bottomrule
  \end{tabular}
\end{table}

\paragraph{Alternative ACDC export.}
For tools requiring YOLO text labels, a separate crowd-free subset contains
349 images and 2,146 boxes. It excludes all 57 images containing crowd regions
and therefore has a different scene distribution. It is a subset of the
406-image validation set, not an independent test set, and its scores must be
reported separately from the full native-JSON evaluation.

\paragraph{Provenance and limitations.}
Exact-file SHA256 comparisons found no matches between ACDC and the checked
foundation/DAWN artifacts, or between MRTMD and the checked
foundation/DAWN/ACDC artifacts. This check does not rule out re-encoded copies,
shared source videos, or perceptual near duplicates. MRTMD contains 919,
387, and 2,405 images from source videos 1, 4, and 5, respectively. These
correlated frames do not constitute 3,711 independent scenes, and only three
video groups are available for assessing between-scene variation. Its 64
bicycle annotations also limit the reliability of bicycle-specific estimates.
Dataset-specific results should therefore be reported separately, with
per-class and per-video breakdowns where practical.

% Reproducibility metadata for supplementary material:
% ACDC images: pcdat1003/ACDC mirror,
% revision e09c38d6e347cd7e8ada42413851d610b311f891.
% ACDC detection ZIP MD5: 32598aacfe0f3c5138262849be8f35f3.
% MRTMD repository: https://github.com/OVD-Labs/MRTMD,
% commit 664f6786e884e41cc95fd2a36949833223acea64.
% MRTMD original annotation SHA256:
% 92d50ba653ccaf362167ef742cee3ea58d42fdbdea28ff351205744d88c8f562.
; requires booktabs in the preamble.
% Bibliography: zero_shot_datasets.bib.
% Draft methods/protocol text, not a report of completed inference experiments.
% Verify all protocol claims against the actual runs before paper submission.
\subsection{Datasets for zero-shot cross-dataset evaluation}
\label{sec:zero-shot-datasets}

We prepare two target datasets for evaluating cross-dataset generalization:
ACDC~\cite{sakaridis2021acdc,sakaridis2026acdc} and
MRTMD~\cite{bugeja2025mrtmd}. Both are represented using the detector's six
classes, in output-index order: person, bicycle, motorcycle, car, bus, and truck.
Here, zero-shot denotes evaluation on a target dataset without target-specific
fitting or hyperparameter selection; it does not denote recognition of unseen
object categories. The proposed evaluation protocol reserves these datasets
for evaluation and fixes preprocessing, checkpoint selection, and detection
thresholds before examining target-set results.

\paragraph{ACDC.}
ACDC comprises driving scenes captured in Switzerland under adverse visual
conditions. We prepare its complete adverse-condition validation split of
406 images at $1920\times1080$ pixels: 100 fog, 106 night, 100 rain, and 100
snow images. Normal-condition reference images are not included.
The source person and rider categories are mapped to person, with bicycle and
motorcycle boxes retained separately. The 46 train annotations are removed
while their images are retained. The resulting six-class annotation set
contains 2,685 individual objects and 63 crowd regions. Native JSON annotations
retain the source boxes, areas, and crowd flags. This mapping defines a custom
six-class evaluation and is not directly comparable with ACDC's official
eight-class detection benchmark.

\paragraph{MRTMD.}
MRTMD comprises traffic-video frames curated from YouTube and distributed at
multiple resolutions. We prepare only the labelled 1080p release, avoiding
repeated counting of different-resolution copies of the same frames.
The pinned source annotations contain 3,733 image records and 69,348 boxes.
Removing 22 redundant image records with identical annotations yields 3,711
images and 69,010 boxes. Class IDs are reordered without merging or dropping
categories. We clip 575 boxes with boundary overflows of at most 0.25 pixels,
and recompute areas from the corrected 1080p boxes because the published area
fields are approximately four times their 1080p bounding-box areas.
The original annotations are retained separately.

\begin{table}[t]
  \centering
  \small
  \caption{Prepared evaluation sets. Crowd annotations describe regions and
  are reported separately from individual object boxes.}
  \label{tab:zero-shot-dataset-sizes}
  \begin{tabular}{lrrr}
    \toprule
    Dataset & Images & Objects & Crowds \\
    \midrule
    ACDC validation & 406 & 2,685 & 63 \\
    MRTMD 1080p & 3,711 & 69,010 & 0 \\
    \bottomrule
  \end{tabular}
\end{table}

\begin{table}[t]
  \centering
  \small
  \caption{Class composition after mapping. ACDC object and crowd counts are
  separate; all MRTMD annotations are individual object boxes.}
  \label{tab:zero-shot-class-counts}
  \begin{tabular}{lrrr}
    \toprule
    Class & \shortstack{ACDC\\objects} & \shortstack{ACDC\\crowds} & MRTMD \\
    \midrule
    Person & 567 & 11 & 21,274 \\
    Bicycle & 68 & 22 & 64 \\
    Motorcycle & 59 & 1 & 15,525 \\
    Car & 1,836 & 29 & 30,187 \\
    Bus & 49 & 0 & 337 \\
    Truck & 106 & 0 & 1,623 \\
    \midrule
    Total & 2,685 & 63 & 69,010 \\
    \bottomrule
  \end{tabular}
\end{table}

\paragraph{Alternative ACDC export.}
For tools requiring YOLO text labels, a separate crowd-free subset contains
349 images and 2,146 boxes. It excludes all 57 images containing crowd regions
and therefore has a different scene distribution. It is a subset of the
406-image validation set, not an independent test set, and its scores must be
reported separately from the full native-JSON evaluation.

\paragraph{Provenance and limitations.}
Exact-file SHA256 comparisons found no matches between ACDC and the checked
foundation/DAWN artifacts, or between MRTMD and the checked
foundation/DAWN/ACDC artifacts. This check does not rule out re-encoded copies,
shared source videos, or perceptual near duplicates. MRTMD contains 919,
387, and 2,405 images from source videos 1, 4, and 5, respectively. These
correlated frames do not constitute 3,711 independent scenes, and only three
video groups are available for assessing between-scene variation. Its 64
bicycle annotations also limit the reliability of bicycle-specific estimates.
Dataset-specific results should therefore be reported separately, with
per-class and per-video breakdowns where practical.

% Reproducibility metadata for supplementary material:
% ACDC images: pcdat1003/ACDC mirror,
% revision e09c38d6e347cd7e8ada42413851d610b311f891.
% ACDC detection ZIP MD5: 32598aacfe0f3c5138262849be8f35f3.
% MRTMD repository: https://github.com/OVD-Labs/MRTMD,
% commit 664f6786e884e41cc95fd2a36949833223acea64.
% MRTMD original annotation SHA256:
% 92d50ba653ccaf362167ef742cee3ea58d42fdbdea28ff351205744d88c8f562.
; requires booktabs in the preamble.
% Bibliography: zero_shot_datasets.bib.
% Draft methods/protocol text, not a report of completed inference experiments.
% Verify all protocol claims against the actual runs before paper submission.
\subsection{Datasets for zero-shot cross-dataset evaluation}
\label{sec:zero-shot-datasets}

We prepare two target datasets for evaluating cross-dataset generalization:
ACDC~\cite{sakaridis2021acdc,sakaridis2026acdc} and
MRTMD~\cite{bugeja2025mrtmd}. Both are represented using the detector's six
classes, in output-index order: person, bicycle, motorcycle, car, bus, and truck.
Here, zero-shot denotes evaluation on a target dataset without target-specific
fitting or hyperparameter selection; it does not denote recognition of unseen
object categories. The proposed evaluation protocol reserves these datasets
for evaluation and fixes preprocessing, checkpoint selection, and detection
thresholds before examining target-set results.

\paragraph{ACDC.}
ACDC comprises driving scenes captured in Switzerland under adverse visual
conditions. We prepare its complete adverse-condition validation split of
406 images at $1920\times1080$ pixels: 100 fog, 106 night, 100 rain, and 100
snow images. Normal-condition reference images are not included.
The source person and rider categories are mapped to person, with bicycle and
motorcycle boxes retained separately. The 46 train annotations are removed
while their images are retained. The resulting six-class annotation set
contains 2,685 individual objects and 63 crowd regions. Native JSON annotations
retain the source boxes, areas, and crowd flags. This mapping defines a custom
six-class evaluation and is not directly comparable with ACDC's official
eight-class detection benchmark.

\paragraph{MRTMD.}
MRTMD comprises traffic-video frames curated from YouTube and distributed at
multiple resolutions. We prepare only the labelled 1080p release, avoiding
repeated counting of different-resolution copies of the same frames.
The pinned source annotations contain 3,733 image records and 69,348 boxes.
Removing 22 redundant image records with identical annotations yields 3,711
images and 69,010 boxes. Class IDs are reordered without merging or dropping
categories. We clip 575 boxes with boundary overflows of at most 0.25 pixels,
and recompute areas from the corrected 1080p boxes because the published area
fields are approximately four times their 1080p bounding-box areas.
The original annotations are retained separately.

\begin{table}[t]
  \centering
  \small
  \caption{Prepared evaluation sets. Crowd annotations describe regions and
  are reported separately from individual object boxes.}
  \label{tab:zero-shot-dataset-sizes}
  \begin{tabular}{lrrr}
    \toprule
    Dataset & Images & Objects & Crowds \\
    \midrule
    ACDC validation & 406 & 2,685 & 63 \\
    MRTMD 1080p & 3,711 & 69,010 & 0 \\
    \bottomrule
  \end{tabular}
\end{table}

\begin{table}[t]
  \centering
  \small
  \caption{Class composition after mapping. ACDC object and crowd counts are
  separate; all MRTMD annotations are individual object boxes.}
  \label{tab:zero-shot-class-counts}
  \begin{tabular}{lrrr}
    \toprule
    Class & \shortstack{ACDC\\objects} & \shortstack{ACDC\\crowds} & MRTMD \\
    \midrule
    Person & 567 & 11 & 21,274 \\
    Bicycle & 68 & 22 & 64 \\
    Motorcycle & 59 & 1 & 15,525 \\
    Car & 1,836 & 29 & 30,187 \\
    Bus & 49 & 0 & 337 \\
    Truck & 106 & 0 & 1,623 \\
    \midrule
    Total & 2,685 & 63 & 69,010 \\
    \bottomrule
  \end{tabular}
\end{table}

\paragraph{Alternative ACDC export.}
For tools requiring YOLO text labels, a separate crowd-free subset contains
349 images and 2,146 boxes. It excludes all 57 images containing crowd regions
and therefore has a different scene distribution. It is a subset of the
406-image validation set, not an independent test set, and its scores must be
reported separately from the full native-JSON evaluation.

\paragraph{Provenance and limitations.}
Exact-file SHA256 comparisons found no matches between ACDC and the checked
foundation/DAWN artifacts, or between MRTMD and the checked
foundation/DAWN/ACDC artifacts. This check does not rule out re-encoded copies,
shared source videos, or perceptual near duplicates. MRTMD contains 919,
387, and 2,405 images from source videos 1, 4, and 5, respectively. These
correlated frames do not constitute 3,711 independent scenes, and only three
video groups are available for assessing between-scene variation. Its 64
bicycle annotations also limit the reliability of bicycle-specific estimates.
Dataset-specific results should therefore be reported separately, with
per-class and per-video breakdowns where practical.

% Reproducibility metadata for supplementary material:
% ACDC images: pcdat1003/ACDC mirror,
% revision e09c38d6e347cd7e8ada42413851d610b311f891.
% ACDC detection ZIP MD5: 32598aacfe0f3c5138262849be8f35f3.
% MRTMD repository: https://github.com/OVD-Labs/MRTMD,
% commit 664f6786e884e41cc95fd2a36949833223acea64.
% MRTMD original annotation SHA256:
% 92d50ba653ccaf362167ef742cee3ea58d42fdbdea28ff351205744d88c8f562.
; requires booktabs in the preamble.
% Bibliography: zero_shot_datasets.bib.
% Draft methods/protocol text, not a report of completed inference experiments.
% Verify all protocol claims against the actual runs before paper submission.
\subsection{Datasets for zero-shot cross-dataset evaluation}
\label{sec:zero-shot-datasets}

We prepare two target datasets for evaluating cross-dataset generalization:
ACDC~\cite{sakaridis2021acdc,sakaridis2026acdc} and
MRTMD~\cite{bugeja2025mrtmd}. Both are represented using the detector's six
classes, in output-index order: person, bicycle, motorcycle, car, bus, and truck.
Here, zero-shot denotes evaluation on a target dataset without target-specific
fitting or hyperparameter selection; it does not denote recognition of unseen
object categories. The proposed evaluation protocol reserves these datasets
for evaluation and fixes preprocessing, checkpoint selection, and detection
thresholds before examining target-set results.

\paragraph{ACDC.}
ACDC comprises driving scenes captured in Switzerland under adverse visual
conditions. We prepare its complete adverse-condition validation split of
406 images at $1920\times1080$ pixels: 100 fog, 106 night, 100 rain, and 100
snow images. Normal-condition reference images are not included.
The source person and rider categories are mapped to person, with bicycle and
motorcycle boxes retained separately. The 46 train annotations are removed
while their images are retained. The resulting six-class annotation set
contains 2,685 individual objects and 63 crowd regions. Native JSON annotations
retain the source boxes, areas, and crowd flags. This mapping defines a custom
six-class evaluation and is not directly comparable with ACDC's official
eight-class detection benchmark.

\paragraph{MRTMD.}
MRTMD comprises traffic-video frames curated from YouTube and distributed at
multiple resolutions. We prepare only the labelled 1080p release, avoiding
repeated counting of different-resolution copies of the same frames.
The pinned source annotations contain 3,733 image records and 69,348 boxes.
Removing 22 redundant image records with identical annotations yields 3,711
images and 69,010 boxes. Class IDs are reordered without merging or dropping
categories. We clip 575 boxes with boundary overflows of at most 0.25 pixels,
and recompute areas from the corrected 1080p boxes because the published area
fields are approximately four times their 1080p bounding-box areas.
The original annotations are retained separately.

\begin{table}[t]
  \centering
  \small
  \caption{Prepared evaluation sets. Crowd annotations describe regions and
  are reported separately from individual object boxes.}
  \label{tab:zero-shot-dataset-sizes}
  \begin{tabular}{lrrr}
    \toprule
    Dataset & Images & Objects & Crowds \\
    \midrule
    ACDC validation & 406 & 2,685 & 63 \\
    MRTMD 1080p & 3,711 & 69,010 & 0 \\
    \bottomrule
  \end{tabular}
\end{table}

\begin{table}[t]
  \centering
  \small
  \caption{Class composition after mapping. ACDC object and crowd counts are
  separate; all MRTMD annotations are individual object boxes.}
  \label{tab:zero-shot-class-counts}
  \begin{tabular}{lrrr}
    \toprule
    Class & \shortstack{ACDC\\objects} & \shortstack{ACDC\\crowds} & MRTMD \\
    \midrule
    Person & 567 & 11 & 21,274 \\
    Bicycle & 68 & 22 & 64 \\
    Motorcycle & 59 & 1 & 15,525 \\
    Car & 1,836 & 29 & 30,187 \\
    Bus & 49 & 0 & 337 \\
    Truck & 106 & 0 & 1,623 \\
    \midrule
    Total & 2,685 & 63 & 69,010 \\
    \bottomrule
  \end{tabular}
\end{table}

\paragraph{Alternative ACDC export.}
For tools requiring YOLO text labels, a separate crowd-free subset contains
349 images and 2,146 boxes. It excludes all 57 images containing crowd regions
and therefore has a different scene distribution. It is a subset of the
406-image validation set, not an independent test set, and its scores must be
reported separately from the full native-JSON evaluation.

\paragraph{Provenance and limitations.}
Exact-file SHA256 comparisons found no matches between ACDC and the checked
foundation/DAWN artifacts, or between MRTMD and the checked
foundation/DAWN/ACDC artifacts. This check does not rule out re-encoded copies,
shared source videos, or perceptual near duplicates. MRTMD contains 919,
387, and 2,405 images from source videos 1, 4, and 5, respectively. These
correlated frames do not constitute 3,711 independent scenes, and only three
video groups are available for assessing between-scene variation. Its 64
bicycle annotations also limit the reliability of bicycle-specific estimates.
Dataset-specific results should therefore be reported separately, with
per-class and per-video breakdowns where practical.

% Reproducibility metadata for supplementary material:
% ACDC images: pcdat1003/ACDC mirror,
% revision e09c38d6e347cd7e8ada42413851d610b311f891.
% ACDC detection ZIP MD5: 32598aacfe0f3c5138262849be8f35f3.
% MRTMD repository: https://github.com/OVD-Labs/MRTMD,
% commit 664f6786e884e41cc95fd2a36949833223acea64.
% MRTMD original annotation SHA256:
% 92d50ba653ccaf362167ef742cee3ea58d42fdbdea28ff351205744d88c8f562.
